\section{Plan-Execute-RePlan 概述}

\subsection{核心理念}

本期介绍了一个非常 general 且在真实科研和工作场景中广泛使用的 Agentic Workflow——\textbf{Dynamic Plan-and-Solve}，即 Plan-Execute-RePlan 三步循环。

\begin{importantbox}{Plan-Execute-RePlan 三步曲}
\begin{enumerate}
    \item \textbf{Plan}：拿到用户的复杂 query/task，先制定一个自顶向下的 sub-goal 分解，输出一个 List of Steps（To-Do List）
    \item \textbf{Execute}：逐步消化 Plan 中的每一个 step，通过工具调用获取 outcome 和 feedback
    \item \textbf{RePlan}：基于执行结果和 feedback，动态更新剩余的 steps
\end{enumerate}
\end{importantbox}

这种模式已被广泛采用——Cursor、Claude Code、Codex 等 Coding Agent 都支持 Plan Mode，先列 To-Do List，再边执行边更新。

\subsection{为什么 Plan-Execute-RePlan 有效}

结构化 Workflow 有助于：
\begin{itemize}
    \item 结构化模型的思考过程
    \item 降低对复杂问题求解的整体复杂度
    \item 自顶向下分解 + 自底向上反馈形成闭环
\end{itemize}

\begin{knowledgebox}{Top-Down Plan vs. Bottom-Up Execute}
\begin{itemize}
    \item \textbf{Plan（Top-Down）}：自顶向下的 sub-goal 分解，不涉及和环境的真实交互，是"刚中之脑"——想象出来的 sub-goal 序列
    \item \textbf{Execute（Bottom-Up）}：和环境真实发生交互，拿到真实的 feedback
    \item \textbf{RePlan}：基于真实交互的结果，动态修正计划
\end{itemize}
\end{knowledgebox}

\subsection{本章小结}

Plan-Execute-RePlan 是一个非常 general 且经过实践验证的 Agentic Workflow，其核心价值在于将复杂问题结构化分解，并通过动态反馈机制弥合规划与执行之间的 gap。

